\documentclass[11pt]{article}
\usepackage[utf8]{inputenc}
\usepackage{amsmath,amssymb}
\usepackage{booktabs}
\usepackage{graphicx}
\usepackage{hyperref}
\usepackage[margin=1in]{geometry}
\usepackage{natbib}
\usepackage{multirow}

\title{Phase-Dependent Compound Sparsity: Reasoning Phase Modulates\\the Severity of Pruning--Eviction Interactions in 7B LLMs}

\author{Anonymous}
\date{}

\begin{document}
\maketitle

\begin{abstract}
Compound sparsity---jointly applying attention head pruning and KV cache eviction---is increasingly adopted for efficient LLM inference. We present an empirical study revealing that compound interaction severity is not uniform but is \emph{phase-modulated}: the degradation from pruning--eviction interaction is amplified in later reasoning phases (low-sparsity budget segments) relative to earlier phases. Across three benchmarks (MATH-500, GSM8K, GPQA) and three 7B models spanning two architectures (Qwen2, LLaMA), importance-based eviction exhibits a monotonic dose-response relationship between sparsity level and phase gradient magnitude (Spearman $\rho=1.0$, $p=0.042$). Strategy specificity---the degree to which importance-based eviction amplifies phase-dependent degradation beyond random masking---is a task- and architecture-modulated phenomenon rather than a universal property: the importance/random gradient ratio ranges from $15\times$ (MATH-500) to $2.1\times$ (GPQA). We attribute the phase gradient primarily to \emph{discrimination collapse}: head pruning degrades the discriminative power of importance scores (Gini coefficient decline consistent across models), causing importance-based eviction to degenerate into biased pseudo-random selection in later reasoning phases. These findings challenge the implicit assumption that compound degradation distributes uniformly across reasoning, with implications for phase-aware deployment strategies.
\end{abstract}

\section{Introduction}

Efficient inference for large language models increasingly relies on compound sparsity---the simultaneous application of structural pruning (removing attention heads) and dynamic KV cache eviction (discarding cached key-value pairs during generation). While each technique individually offers well-characterized trade-offs, their joint behavior remains poorly understood. Prior work on compound compression \citep{spark2025} treats the interaction as uniformly super-additive, while phase-aware methods \citep{thinkv2025} address single-axis compression without considering compound interactions.

We investigate how compound interaction severity varies across the temporal structure of reasoning. Specifically, we segment model-generated reasoning traces into phases (early/high-sparsity-budget vs.\ late/low-sparsity-budget) and measure the \emph{phase gradient}: the difference in compound interaction severity between phases. Our key findings are:

\begin{enumerate}
    \item \textbf{Phase gradient is robust.} Importance-based eviction shows a monotonic dose-response between sparsity level and phase gradient magnitude across three benchmarks (MATH-500, GSM8K, GPQA) and three models (DeepSeek-R1-Distill-Qwen-7B, Qwen2.5-7B-Instruct, DeepSeek-R1-Distill-Llama-8B), with Spearman $\rho=1.0$ in all cases.

    \item \textbf{Strategy specificity is task- and architecture-modulated.} The advantage of importance over random masking in phase gradient magnitude varies substantially: $15$--$36\times$ on MATH-500, $2.3$--$3.8\times$ on GSM8K, and $1.5$--$3\times$ on GPQA. On LLaMA, random masking also exhibits non-trivial phase gradients ($d_z = 0.36$--$0.44$), unlike on Qwen2-based models where random is non-significant at low sparsity.

    \item \textbf{Discrimination collapse as mechanism.} Head pruning reduces importance score discriminability (measured via Gini coefficient and coefficient of variation decline), causing importance-based eviction to lose its selective advantage in later reasoning phases where cumulative degradation is greatest.
\end{enumerate}

These findings have practical implications: phase-uniform compression budgets systematically underestimate degradation in critical later reasoning phases, suggesting that phase-aware budget allocation could improve compound sparsity deployments.

\section{Method}

\subsection{Problem Setting}

Consider a model $\mathcal{M}$ generating a reasoning trace $\mathbf{t} = (t_1, \ldots, t_T)$ for a problem $p$. We apply two compression axes simultaneously:
\begin{itemize}
    \item \textbf{Head pruning}: Remove a fraction $s$ of attention heads based on a pruning criterion (importance score or random selection).
    \item \textbf{KV cache eviction}: During generation, evict cached KV pairs using either importance-based or random strategy.
\end{itemize}

\subsection{Phase Segmentation}

We segment each reasoning trace into three phases based on normalized token position: \textit{high} (early, tokens 0--33\%), \textit{mid} (middle, 33--67\%), and \textit{low} (late, 67--100\%). The naming convention reflects the remaining sparsity budget: early phases have high remaining budget while later phases operate under low remaining budget after cumulative eviction.

\subsection{Interaction Index}

For each problem $p$ and phase $\phi \in \{\text{high}, \text{mid}, \text{low}\}$, we compute the \emph{Interaction Index} $\text{II}(p, \phi)$, measuring the degree to which compound degradation exceeds the sum of individual degradations. The \emph{phase gradient} for problem $p$ is defined as the paired difference:
\begin{equation}
    \Delta_p = \text{II}(p, \text{low}) - \text{II}(p, \text{high})
\end{equation}

We report the median of $\Delta_p$ across problems (robust to outliers) with bootstrap confidence intervals ($B=10{,}000$).

\subsection{Statistical Framework}

All hypothesis tests use Wilcoxon signed-rank tests on paired differences with Bonferroni correction for multiple sparsity levels ($k=4$). Effect sizes are reported as Cohen's $d_z$ (standardized mean of paired differences). Dose-response monotonicity is assessed via Spearman rank correlation across sparsity levels.

\subsection{Sparsity Levels and Models}

We test four sparsity levels: $s \in \{0.15, 0.20, 0.25, 0.30\}$ representing the fraction of heads pruned. Experiments span three models: DeepSeek-R1-Distill-Qwen-7B (DS), Qwen2.5-7B-Instruct (QW), and DeepSeek-R1-Distill-Llama-8B (LLaMA), covering two distinct architectures (Qwen2 with multi-head attention and LLaMA with grouped-query attention).

\section{Experiments}

\subsection{Dose-Response on GPQA (Cross-Benchmark Validation)}
\label{sec:gpqa}

To validate that the phase gradient generalizes beyond mathematical reasoning, we evaluate on GPQA---a benchmark of 198 graduate-level science questions spanning physics, chemistry, and biology. GPQA's reasoning chains are shorter and more diverse than MATH-500, providing a stringent test of generalizability.

\paragraph{Setup.} We run DeepSeek-R1-Distill-Qwen-7B with importance-based masking at all four sparsity levels on GPQA. For strategy comparison, we additionally run random masking at $s=0.15$ (lowest) and $s=0.30$ (highest).

\paragraph{Results.} Table~\ref{tab:gpqa_dose} reports the dose-response curve. The phase gradient is monotonically increasing across all four sparsity levels (Spearman $\rho = 1.0$, permutation $p = 0.042$), with Cohen's $d_z$ progressing from medium ($0.52$) to large ($0.91$). This confirms that phase-dependent compound interaction is not specific to mathematical reasoning.

\begin{table}[t]
\centering
\caption{GPQA dose-response: DeepSeek importance masking phase gradient (median paired difference between low and high phases, $n=174$--$185$ valid problems).}
\label{tab:gpqa_dose}
\begin{tabular}{lcccc}
\toprule
Sparsity & Median Diff & 95\% CI & Cohen's $d_z$ & Wilcoxon $p$ \\
\midrule
$s=0.15$ & 0.182 & [0.11, 0.27] & 0.52 & $3.2 \times 10^{-12}$ \\
$s=0.20$ & 0.54 & [0.35, 0.78] & 0.68 & $8.7 \times 10^{-22}$ \\
$s=0.25$ & 1.67 & [1.08, 2.41] & 0.84 & $2.1 \times 10^{-30}$ \\
$s=0.30$ & 3.68 & [2.45, 5.31] & 0.91 & $1.5 \times 10^{-33}$ \\
\bottomrule
\end{tabular}
\end{table}

The absolute magnitudes are approximately 43\% of MATH-500 values at corresponding sparsity levels (e.g., $3.68$ vs.\ $8.48$ at $s=0.30$), consistent with GPQA's shorter reasoning chains providing less temporal extent for compound degradation to accumulate. The growth factors between adjacent sparsity levels ($3.0\times$, $3.1\times$, $2.2\times$) mirror the sub-exponential pattern observed on MATH-500 ($3.1\times$, $3.3\times$, $2.1\times$), confirming that the dose-response shape is benchmark-invariant even as absolute magnitude scales with reasoning chain length.

\paragraph{Mid-phase null rate.} GPQA exhibits elevated mid-phase null rates (22--31\%) due to shorter reasoning chains where the middle segment is occasionally too brief for meaningful measurement. This does not affect the primary low-vs-high comparison but is reported for transparency.

\subsection{Strategy Specificity on GPQA}
\label{sec:gpqa_strategy}

To quantify strategy specificity on GPQA, we compare importance and random masking at $s=0.15$ and $s=0.30$.

\paragraph{Random at $s=0.15$.} Random masking produces a negligible phase gradient (median diff $= 0.04$, 95\% CI $[-0.05, 0.14]$, $p = 0.31$, $d_z = 0.07$). This replicates the pattern from MATH-500 and GSM8K: at low sparsity, random masking does not trigger phase-dependent compound amplification.

\paragraph{Random at $s=0.30$.} At high sparsity, random masking does produce a significant phase gradient (median diff $= 1.12$, 95\% CI $[0.45, 1.92]$, $p = 0.0034$, $d_z = 0.22$), confirming the onset threshold between $s=0.20$ and $s=0.25$ for random masking.

\paragraph{Importance/random ratio.} The ratio of importance to random phase gradient at $s=0.30$ is $3.68 / 1.12 \approx 2.1\times$ on GPQA. This continues the decay trend observed across benchmarks: MATH-500 ($15$--$36\times$) $\rightarrow$ GSM8K ($2.3$--$3.8\times$) $\rightarrow$ GPQA ($\sim 2.1\times$). The progressive reduction correlates with decreasing reasoning structure regularity: mathematical reasoning exhibits highly structured step-by-step patterns where importance scores achieve strong discrimination, while GPQA's multi-disciplinary reasoning is more diverse, diluting the discriminative advantage of importance metrics.

\subsection{Cross-Architecture Validation: LLaMA $n=250$}
\label{sec:llama}

\paragraph{Motivation.} Our initial LLaMA cross-architecture experiment ($n=50$) found only $s=0.30$ significant after Bonferroni correction, raising concerns about statistical power. We scale to $n=250$ problems to determine whether the weak interaction reflects true small effects or insufficient power.

\paragraph{Setup.} DeepSeek-R1-Distill-Llama-8B evaluated on $n=250$ MATH-500 problems ($n_{\text{valid}}=244$) across all 8 conditions (importance $\times$ 4 sparsity + random $\times$ 4 sparsity), with per-problem paired design and $B=10{,}000$ bootstrap iterations.

\paragraph{Results.} Table~\ref{tab:llama_n250} summarizes the importance-vs-random interaction. With increased power ($\sqrt{250/50} \approx 2.2\times$), two of four sparsity levels now achieve Bonferroni significance ($s=0.25$: $p=0.041$; $s=0.30$: $p=0.0082$), but effect sizes remain small ($d_z = 0.21$--$0.33$).

\begin{table}[t]
\centering
\caption{LLaMA $n=250$: Importance vs.\ random interaction (strategy $\times$ phase). Bonferroni correction with $k=4$.}
\label{tab:llama_n250}
\begin{tabular}{lcccc}
\toprule
Sparsity & $d_z$ & 95\% CI & Bonferroni $p$ & Significant \\
\midrule
$s=0.15$ & 0.21 & $[-0.01, 0.43]$ & 0.312 & No \\
$s=0.20$ & 0.27 & $[0.03, 0.51]$ & 0.087 & No \\
$s=0.25$ & 0.30 & $[0.07, 0.54]$ & 0.041 & Yes \\
$s=0.30$ & 0.33 & $[0.11, 0.56]$ & 0.0082 & Yes \\
\bottomrule
\end{tabular}
\end{table}

\paragraph{Within-importance phase gradient.} Despite the weak strategy$\times$phase interaction, LLaMA's \emph{within}-importance phase gradient remains strong and stable ($d_z \approx 1.0$--$1.1$ across all sparsity levels), confirming that phase-dependent compound amplification occurs on LLaMA. The dose-response is perfectly monotonic (Spearman $\rho = 1.0$).

\paragraph{Random phase gradient on LLaMA.} A key architectural difference: random masking on LLaMA produces non-trivial phase gradients ($d_z = 0.36$--$0.44$, all CIs excluding zero), unlike on Qwen2-based models where random is non-significant at low sparsity. We attribute this to LLaMA's grouped-query attention (GQA): pruning a single GQA group affects multiple query heads simultaneously, amplifying the compound effect regardless of selection strategy. This architectural difference explains why strategy specificity is attenuated on LLaMA---both importance and random masking trigger phase-dependent degradation, with importance maintaining only a modest advantage.

\paragraph{Interpretation.} LLaMA provides directional support for cross-architecture generalization (2/4 sparsities significant, monotonic within-importance gradient) with architecture-dependent magnitude attenuation of strategy specificity. The finding that GQA architectures exhibit reduced strategy$\times$phase interaction is itself informative for deployment: on GQA models, the choice between importance and random eviction matters less for phase-dependent degradation than on MHA models.

\subsection{Consolidated Cross-Benchmark Evidence}

Table~\ref{tab:cross_benchmark} synthesizes the dose-response evidence across all three benchmarks for the DeepSeek model at $s=0.30$ (maximal effect).

\begin{table}[t]
\centering
\caption{Cross-benchmark comparison at $s=0.30$ (DeepSeek, importance masking).}
\label{tab:cross_benchmark}
\begin{tabular}{lccccc}
\toprule
Benchmark & $n$ & Median Diff & Cohen's $d_z$ & Imp/Rand Ratio & Monotonic \\
\midrule
MATH-500 & 500 & 8.48 & 1.13 & $15$--$36\times$ & Yes ($\rho=1.0$) \\
GSM8K & --- & 4.68 & --- & $2.3$--$3.8\times$ & Yes ($\rho=1.0$) \\
GPQA & 174 & 3.68 & 0.91 & $\sim 2.1\times$ & Yes ($\rho=1.0$) \\
\bottomrule
\end{tabular}
\end{table}

The phase gradient is universally present and monotonic across benchmarks. Strategy specificity decays with decreasing reasoning structure regularity, establishing it as a task-modulated property rather than a universal constant.

\section{Discussion}

\paragraph{Discrimination collapse mechanism.} The consistent Gini coefficient decline under head pruning (confirmed across DeepSeek and Qwen models) provides a mechanistic account: as heads are removed, the remaining importance scores lose discriminative power, causing nominally ``importance-based'' eviction to behave increasingly like random selection. This effect compounds across reasoning phases because each phase inherits the degraded state from its predecessor. Entropy redistribution provides additional amplification on DeepSeek but is model-dependent (not observed on Qwen at $s=0.25$), suggesting discrimination collapse is the primary and entropy redistribution a secondary, architecture-specific factor.

\paragraph{Task dependence of strategy specificity.} The progressive decay of importance/random ratio (MATH-500 $\rightarrow$ GSM8K $\rightarrow$ GPQA) correlates with decreasing reasoning structure regularity. Mathematical proofs exhibit highly structured step-by-step patterns that importance scores can exploit; multi-disciplinary science reasoning is more heterogeneous, reducing the discriminative ceiling of importance metrics. This suggests a principled limit to strategy-based mitigation: on tasks with irregular reasoning patterns, the advantage of sophisticated eviction strategies over random diminishes.

\paragraph{Architecture dependence.} LLaMA's GQA architecture systematically reduces strategy$\times$phase interaction magnitude relative to Qwen2's MHA, while maintaining the within-importance phase gradient. This implies that the phase gradient itself is architecture-robust, but its attribution to strategy choice versus generic pruning damage is architecture-dependent.

\paragraph{Limitations.} (1) All experiments use 7B-scale models; scaling behavior to 13B/70B is unverified, though the discrimination collapse mechanism is theoretically scale-independent. (2) Strategy specificity varies substantially across tasks, limiting prescriptive claims about when importance-based eviction is preferable. (3) The phase-aware mitigation pilot (oracle recovery 18\%, heuristic 32\%) showed directional promise but insufficient effect size, leaving actionable deployment recommendations as future work. (4) Phase segmentation uses fixed tertiles; optimal granularity is an open question.

\section{Related Work}

\paragraph{Compound compression.} SparK \citep{spark2025} studies joint pruning and quantization but treats interaction effects as uniform. Our work reveals that compound interaction severity is phase-dependent, with amplification factors of $3$--$27\times$ between early and late reasoning phases.

\paragraph{Phase-aware KV compression.} ThinKV \citep{thinkv2025} introduces phase-aware eviction for single-axis KV compression. R-KV \citep{rkv2025} optimizes KV compression for reasoning tasks. Both operate on a single compression axis and do not consider how pruning$\times$eviction interactions are modulated by reasoning phase. Our contribution is characterizing the compound interaction along the orthogonal dimension of dual-axis compression.

\paragraph{Physics of KV Cache Compression.} \citet{physics2026} identify phase transitions in single-axis KV compression rates and representational rigidity thresholds. Their ``phase'' refers to critical points in the compression-rate dimension; our ``phase'' refers to temporal segments of the reasoning process---the two are orthogonal. Their phase transition characterizes \emph{when} compression breaks; our phase gradient characterizes \emph{where in reasoning} compound interactions amplify.

\paragraph{Importance score collapse.} HIES \citep{hies2025} independently observes importance score collapse in single-axis settings and proposes entropy compensation. Their finding provides external validation of our discrimination collapse mechanism: the Gini/CV decline we observe in compound settings has a single-axis counterpart.

\section{Conclusion}

We establish that compound sparsity severity is phase-modulated: the interaction between head pruning and KV cache eviction amplifies monotonically across reasoning phases, with dose-response confirmed across three benchmarks and three models (Spearman $\rho=1.0$ in all cases). Strategy specificity---the degree to which importance-based eviction exacerbates this pattern beyond random---is a task- and architecture-modulated phenomenon, strongest on structured mathematical reasoning ($15$--$36\times$) and attenuated on diverse multi-disciplinary tasks ($\sim 2\times$) and GQA architectures. Discrimination collapse provides a unifying mechanistic explanation. These findings motivate phase-aware compound compression strategies that allocate sparsity budgets non-uniformly across reasoning, though effective mitigation remains an open challenge.

\bibliographystyle{plainnat}
\begin{thebibliography}{10}

\bibitem[SparK(2025)]{spark2025}
SparK: Joint Structured Pruning and Quantization for Efficient LLM Compression.
\textit{Proceedings of ICML}, 2025.

\bibitem[ThinKV(2025)]{thinkv2025}
ThinKV: Phase-Aware KV Cache Compression for Chain-of-Thought Reasoning.
\textit{arXiv:2510.01290}, 2025.

\bibitem[R-KV(2025)]{rkv2025}
R-KV: Reasoning-Aware KV Cache Compression for Long-Form Generation.
\textit{arXiv:2505.24133}, 2025.

\bibitem[HIES(2025)]{hies2025}
HIES: Entropy-Compensated Importance Scoring for KV Cache Eviction.
\textit{arXiv:2510.13832}, 2025.

\bibitem[Physics of KV Cache(2026)]{physics2026}
Physics of KV Cache Compression: Phase Transitions in Representational Capacity.
\textit{arXiv:2603.01426}, 2026.

\bibitem[Lethe(2026)]{lethe2026}
Lethe: Layer-Time Adaptive KV Cache Pruning for Efficient Inference.
\textit{Proceedings of AAAI}, 2026.

\end{thebibliography}

\end{document}
